% GENERATED FILE. Edit canonical lessons, then run npm run generate:books.
% canonical-source-hash: fnv1a64:5fceb3aa4c077cea
% canonical-lessons: GU-C482-kakhghodi, GU-C482-dhimdu, GU-C482-phulvu, GU-C482-cattu, GU-C482-aphalvu

\chapter{A Bad Fall}
\label{ch:gu-a-bad-fall}
\hlchaptermodality{482}

\begin{chapteropening}
\textbf{By the end of this chapter:} \emph{I can say how I fell and what hurt.}

Complete the last lesson of chapter 482: I can say how I fell and what hurt.
\end{chapteropening}

\section[\texorpdfstring{\gu{કાખઘોડી} (kākhghoḍī)}{કાખઘોડી (kākhghoḍī)}]{\gu{કાખઘોડી} (kākhghoḍī) — a crutch}

\label{lesson:GU-C482-kakhghodi}



\begin{quote}
Before the new one: say the Gujarati for to ooze, then the Gujarati for to get grazed.
\end{quote}

\subsection*{You'll want to know: \gu{કાખઘોડી}}
\textbf{\gu{કાખઘોડી}} — \emph{kākhghoḍī} — \textquotedblleft{}a crutch\textquotedblright{}.

\begin{quote}
\gu{પડ્યા} \gu{પછી} \gu{દાદા} \gu{કાખઘોડીથી} \gu{ચાલે} \gu{છે}.
\end{quote}

\begin{quote}
\emph{paḍyā pachī dādā kākhghoḍīthī cāle che.}
\end{quote}

\begin{quote}
Since his fall my grandfather walks with a crutch.
\end{quote}

A feminine noun.

\begin{cousinweb}
\textbf{\gu{કાખઘોડી}} (\emph{kākhghoḍī}) is built from \emph{kākh}, \textquotedblleft{}armpit\textquotedblright{}, which CDIAL lists under Sanskrit \emph{kakṣa} (CDIAL 2588), and \emph{ghoḍī}, \textquotedblleft{}a little horse\textquotedblright{}, from the family of \textbf{\gu{ઘોડો}} (\emph{ghoḍo}) (CDIAL 4516): a little horse under the arm.
\end{cousinweb}

\begin{grammarlens}[title={What you've built}]
One more word for describing what you lived and felt.
\end{grammarlens}

\subsection*{Your turn}
\begin{itemize}
\raggedright
  \item \emph{Say it:} \emph{kākhghoḍī}
  \item \emph{Say it:} \emph{kākhghoḍī}, once more
  \item \emph{Say it:} \emph{paḍyā pachī dādā kākhghoḍīthī cāle che.} — Since his fall my grandfather walks with a crutch.
  \item \emph{Recall:} say the Gujarati for to ooze, then the Gujarati for to get grazed, then say \emph{kākhghoḍī} again
\end{itemize}

\begin{tcolorbox}[breakable,colback=teal!4,colframe=teal!35!black,title={Before you move on}]
What does \gu{કાખઘોડી} mean? (\textquotedblleft{}A crutch\textquotedblright{}.) Say it once more.
\end{tcolorbox}
\section[\texorpdfstring{\gu{ઢીમડું} (ḍhīmḍũ)}{ઢીમડું (ḍhīmḍũ)}]{\gu{ઢીમડું} (ḍhīmḍũ) — a bump (from a knock)}

\label{lesson:GU-C482-dhimdu}



\begin{quote}
Before the new one: say the Gujarati for to get grazed, then the Gujarati for a crutch.
\end{quote}

\subsection*{You'll want to know: \gu{ઢીમડું}}
\textbf{\gu{ઢીમડું}} — \emph{ḍhīmḍũ} — \textquotedblleft{}a bump (from a knock)\textquotedblright{}.

\begin{quote}
\gu{મારા} \gu{માથા} \gu{પર} \gu{મોટું} \gu{ઢીમડું} \gu{થયું}.
\end{quote}

\begin{quote}
\emph{mārā māthā par moṭũ ḍhīmḍũ thayũ.}
\end{quote}

\begin{quote}
A big bump came up on my head.
\end{quote}

A neuter noun.

\begin{cousinweb}
\textbf{\gu{ઢીમડું}} (\emph{ḍhīmḍũ}) goes back to a reconstructed \emph{*ḍhīmma-}, \textquotedblleft{}lump\textquotedblright{}; CDIAL lists the Gujarati word, \textquotedblleft{}big swelling or bump\textquotedblright{} (CDIAL 5591).
\end{cousinweb}

\begin{grammarlens}[title={What you've built}]
One more word for describing what you lived and felt.
\end{grammarlens}

\subsection*{Your turn}
\begin{itemize}
\raggedright
  \item \emph{Say it:} \emph{ḍhīmḍũ}
  \item \emph{Say it:} \emph{ḍhīmḍũ}, once more
  \item \emph{Say it:} \emph{mārā māthā par moṭũ ḍhīmḍũ thayũ.} — A big bump came up on my head.
  \item \emph{Recall:} say the Gujarati for to get grazed, then the Gujarati for a crutch, then say \emph{ḍhīmḍũ} again
\end{itemize}

\begin{tcolorbox}[breakable,colback=teal!4,colframe=teal!35!black,title={Before you move on}]
What does \gu{ઢીમડું} mean? (\textquotedblleft{}A bump (from a knock)\textquotedblright{}.) Say it once more.
\end{tcolorbox}
\section[\texorpdfstring{\gu{ફૂલવું} (phūlvũ)}{ફૂલવું (phūlvũ)}]{\gu{ફૂલવું} (phūlvũ) — to swell up; to puff out}

\label{lesson:GU-C482-phulvu}



\begin{quote}
Before the new one: say the Gujarati for a crutch, then the Gujarati for a bump.
\end{quote}

\subsection*{You'll want to know: \gu{ફૂલવું}}
\textbf{\gu{ફૂલવું}} — \emph{phūlvũ} — \textquotedblleft{}to swell up; to puff out\textquotedblright{}.

\begin{quote}
\gu{પડ્યા} \gu{પછી} \gu{પગની} \gu{ઘૂંટી} \gu{ફૂલી} \gu{ગઈ}.
\end{quote}

\begin{quote}
\emph{paḍyā pachī pagnī ghū̃ṭī phūlī gaī.}
\end{quote}

\begin{quote}
After the fall my ankle swelled up.
\end{quote}

Past (it …): \gu{તે} \gu{ઘૂંટી} \gu{ફૂલી} \gu{ગઈ}, \emph{te ghū̃ṭī phūlī gaī}.

\begin{cousinweb}
\textbf{\gu{ફૂલવું}} (\emph{phūlvũ}) is inherited from Sanskrit \emph{phullati}, \textquotedblleft{}expands, opens like a flower\textquotedblright{}; CDIAL lists the Gujarati verb, \textquotedblleft{}to swell, blossom\textquotedblright{} (CDIAL 9093).
\end{cousinweb}

\begin{grammarlens}[title={What you've built}]
One more word for describing what you lived and felt.
\end{grammarlens}

\subsection*{Your turn}
\begin{itemize}
\raggedright
  \item \emph{Say it:} \emph{phūlvũ}
  \item \emph{Say it:} \emph{phūlvũ}, once more
  \item \emph{Say it:} \emph{paḍyā pachī pagnī ghū̃ṭī phūlī gaī.} — After the fall my ankle swelled up.
  \item \emph{Recall:} say the Gujarati for a crutch, then the Gujarati for a bump, then say \emph{phūlvũ} again
\end{itemize}

\begin{tcolorbox}[breakable,colback=teal!4,colframe=teal!35!black,title={Before you move on}]
What does \gu{ફૂલવું} mean? (\textquotedblleft{}To swell up; to puff out\textquotedblright{}.) Say it once more.
\end{tcolorbox}
\section[\texorpdfstring{\gu{ચત્તું} (cattũ)}{ચત્તું (cattũ)}]{\gu{ચત્તું} (cattũ) — flat on one's back}

\label{lesson:GU-C482-cattu}



\begin{quote}
Before the new one: say the Gujarati for a bump, then the Gujarati for to swell up.
\end{quote}

\subsection*{You'll want to know: \gu{ચત્તું}}
\textbf{\gu{ચત્તું}} — \emph{cattũ} — \textquotedblleft{}flat on one's back\textquotedblright{}.

\begin{quote}
\gu{તે} \gu{ભીની} \gu{ભોંય} \gu{પર} \gu{ચત્તો} \gu{પડ્યો}.
\end{quote}

\begin{quote}
\emph{te bhīnī bhõy par catto paḍyo.}
\end{quote}

\begin{quote}
He fell flat on his back on the wet floor.
\end{quote}

\begin{cousinweb}
\textbf{\gu{ચત્તું}} (\emph{cattũ}) goes back to a reconstructed \emph{*citta-}, \textquotedblleft{}lying on the back\textquotedblright{}; CDIAL lists the Gujarati word (CDIAL 4800).
\end{cousinweb}

\begin{grammarlens}[title={What you've built}]
One more word for describing what you lived and felt.
\end{grammarlens}

\subsection*{Your turn}
\begin{itemize}
\raggedright
  \item \emph{Say it:} \emph{cattũ}
  \item \emph{Say it:} \emph{cattũ}, once more
  \item \emph{Say it:} \emph{te bhīnī bhõy par catto paḍyo.} — He fell flat on his back on the wet floor.
  \item \emph{Recall:} say the Gujarati for a bump, then the Gujarati for to swell up, then say \emph{cattũ} again
\end{itemize}

\begin{tcolorbox}[breakable,colback=teal!4,colframe=teal!35!black,title={Before you move on}]
What does \gu{ચત્તું} mean? (\textquotedblleft{}Flat on one's back\textquotedblright{}.) Say it once more.
\end{tcolorbox}
\section[\texorpdfstring{\gu{અફાળવું} (aphāḷvũ)}{અફાળવું (aphāḷvũ)}]{\gu{અફાળવું} (aphāḷvũ) — to bang; to dash against}

\label{lesson:GU-C482-aphalvu}



\begin{quote}
Before the new one: say the Gujarati for to swell up, then the Gujarati for flat on one's back.
\end{quote}

\subsection*{You'll want to know: \gu{અફાળવું}}
\textbf{\gu{અફાળવું}} — \emph{aphāḷvũ} — \textquotedblleft{}to bang; to dash against\textquotedblright{}.

\begin{quote}
\gu{અંધારામાં} \gu{મેં} \gu{માથું} \gu{દીવાલ} \gu{સાથે} \gu{અફાળ્યું}.
\end{quote}

\begin{quote}
\emph{andhārāmā̃ mẽ māthũ dīvāl sāthe aphāḷyũ.}
\end{quote}

\begin{quote}
In the dark I banged my head on the wall.
\end{quote}

Past, with \gu{તેણે} (he or she …): \gu{તેણે} \gu{માથું} \gu{અફાળ્યું}, \emph{teṇe māthũ aphāḷyũ}.

\begin{cousinweb}
\textbf{\gu{અફાળવું}} (\emph{aphāḷvũ}) is inherited from Sanskrit \emph{āsphālayati}, \textquotedblleft{}causes to flap, shakes, throws\textquotedblright{}; CDIAL lists the Gujarati verb, \textquotedblleft{}to dash against\textquotedblright{} (CDIAL 1531).
\end{cousinweb}

\begin{grammarlens}[title={What you've built}]
One more word for describing what you lived and felt.
\end{grammarlens}

\subsection*{Your turn}
\begin{itemize}
\raggedright
  \item \emph{Say it:} \emph{aphāḷvũ}
  \item \emph{Say it:} \emph{aphāḷvũ}, once more
  \item \emph{Say it:} \emph{andhārāmā̃ mẽ māthũ dīvāl sāthe aphāḷyũ.} — In the dark I banged my head on the wall.
  \item \emph{Recall:} say the Gujarati for to swell up, then the Gujarati for flat on one's back, then say \emph{aphāḷvũ} again
\end{itemize}

\begin{tcolorbox}[breakable,colback=teal!4,colframe=teal!35!black,title={Before you move on}]
What does \gu{અફાળવું} mean? (\textquotedblleft{}To bang; to dash against\textquotedblright{}.) Say it once more.
\end{tcolorbox}
